% Generated draft: requires mathematical and source review.
\documentclass[11pt]{article}
\usepackage[margin=25mm]{geometry}
\usepackage{fontspec}
\setmainfont{Latin Modern Roman}
\usepackage{amsmath,amssymb,mathtools,unicode-math}
\setmathfont{Latin Modern Math}
\usepackage{enumitem,needspace,xurl,hyperref}
\setlength{\parindent}{0pt}
\setlength{\parskip}{5pt}
\setlength{\emergencystretch}{3em}
\newcommand{\N}{\mathbb{N}}
\newcommand{\Z}{\mathbb{Z}}
\newcommand{\Q}{\mathbb{Q}}
\newcommand{\R}{\mathbb{R}}
\newcommand{\C}{\mathbb{C}}
\newcommand{\F}{\mathbb{F}}
\newcommand{\A}{\mathbb{A}}
\newcommand{\PP}{\mathbb{P}}
\newcommand{\E}{\mathbb{E}}
\newcommand{\eps}{\varepsilon}
\newcommand{\dd}{\,\mathrm{d}}
\begin{document}
\section*{UnsolvedMath ID: 30003854}
\textbf{Source problem code:} catalogue entry
\subsection{Definitions and mathematical statement}
For a group \(H\), a classifying space \(K(H,1)\) is a connected CW complex with fundamental group \(H\) and contractible universal cover. The group \(H\) is of type \(F_\infty\) if it admits a classifying space whose \(n\)-skeleton is finite for every integer \(n\geq 0\). Here the source notation \(F^{\infty}\) is interpreted as the standard finiteness property \(F_\infty\); this interpretation requires review.

The group \(H\) is of type \(FP_2\) over \(\Z\) if the trivial left \(\Z H\)-module \(\Z\) has a projective resolution
\[
\cdots\longrightarrow P_2\longrightarrow P_1\longrightarrow P_0\longrightarrow\Z\longrightarrow 0
\]
in which \(P_0,P_1,P_2\) are finitely generated. The coefficient ring is taken to be \(\Z\), following the usual convention for an unqualified \(FP_2\); the supplied statement does not specify coefficients. A group is finitely presentable if it admits a presentation with finitely many generators and finitely many defining relations.

The supplied record combines two questions, which must be kept separate.
\begin{enumerate}
\item For the intended Lodha--Moore groups \(G\), equipped with the source's designated natural homomorphism \(\chi_G:G\to\Z\), is
\[
\ker(\chi_G)=\{g\in G:\chi_G(g)=0\}
\]
of type \(F_\infty\)? This is only a schematic formulation: the supplied record neither defines the intended groups nor specifies \(\chi_G\). It also does not explicitly state that \(\chi_G\) is surjective. Thus neither the precise range of \(G\) nor the kernel is determined by the supplied statement.
\item Let \(\mathcal T\) denote the intended family of Thompson-like groups. Does there exist a group \(H\in\mathcal T\) such that \(H\) is of type \(FP_2\) over \(\Z\) and is not finitely presentable? The supplied record gives no definition of \(\mathcal T\), so this existential question has no fully specified domain here.
\end{enumerate}

The supplied verification note states that the extraction merges two independent numbered questions in the report. Their separate original numbers and complete formulations are not supplied. Consequently these are qualified reconstructions, not a recovered single precise problem. The registry marks the statement as \texttt{unrecoverable} and the overall literature status as \texttt{partially\_solved}; those are supplied curation labels.
\subsection{Short English statement}
The record combines two independent questions: whether certain unspecified Lodha--Moore kernels have finite skeleta in every dimension in a classifying space, and whether the intended family of Thompson-like groups contains a group of type \(FP_2\) that is not finitely presentable. The groups, natural homomorphisms, and intended family need specification. Supplied dated curation reports the first question resolved and the second open; this draft does not independently verify that assessment.
\subsection{Sources}
\begin{itemize}
\item Supplied local UnsolvedMath registry record, permanent ID \(30003854\), source code \texttt{catalogue entry}, title ``Finiteness Properties of Lodha--Moore Kernels.'' No individual proposer is supplied. The proposed year is recorded as 2018.
\item Supplied report citation: \emph{Cohomological and Metric Properties of Groups of Homeomorphisms of R} (2019), Oberwolfach Reports, pp.~1579--1633, DOI \texttt{10.4171/owr/2018/26}. \url{https://doi.org/10.4171/owr/2018/26}. The supplied citation uses 2019, while the proposed year and report identifier refer to 2018; these source details are retained without reconciliation.
\item Lodha--Zaremsky, \emph{The BNSR invariants of the Lodha--Moore groups}: \url{https://arxiv.org/abs/2007.12518}.
\item Supplied Cambridge University Press article link: \url{https://www.cambridge.org/core/services/aop-cambridge-core/content/view/04297B3A2B295AB357EA34BA136473C5/S0305004122000111a.pdf}.
\item Supplied Oberwolfach Report 26/2018 link: \url{https://ems.press/content/serial-article-files/46748?nt=1}.
\item Literature assessment dated exactly \texttt{2026-08-21}: supplied dated curation, not independently verified. It attributes resolution of the Lodha--Moore kernel question to Lodha--Zaremsky through determination of higher BNSR invariants and a proof that the relevant finitely presented normal subgroups are of type \(F_\infty\). It reports that the separately appended \(FP_2\) question remains open. The background explicitly calls this dated literature triage derived from the supplied OWR curation and requests independent verification of time-sensitive status and cited claims.
\item Supplied statement-verification note: the official report is reported to contain two independent numbered questions merged by this extraction. No separate clean statement was promoted, and the original record was retained publicly for provenance. This source-verification assertion has not been independently checked for this draft.
\end{itemize}
\end{document}
